\documentclass[11pt]{article}
\usepackage{mystyle}

\title{FIT1058 --- Chapter 2: Functions\\ \large Exercise Solutions}
\author{Alston}
\date{Semester 2, 2026}

\begin{document}
\maketitle

% ======================================================================
\begin{exercise}{1}
    Let $A$ be a finite set. How many indicator functions with domain
    $A$ are there?
\end{exercise}
\begin{solution}
    An indicator function is defined by
    \[
        \chi_A(x) = \begin{cases} 1 & \text{if } x\in A \\
            0 & \text{if } x\notin A \end{cases}
    \]

    Let $|A|=n$. The number of functions $\chi: A\to\{0,1\}$ equals
    the number of ways to independently assign each of the $n$
    elements of $A$ a value in $\{0,1\}$, which is $2^n$.

    Equivalently, each such function corresponds to exactly one
    subset $X\subseteq A$ (namely $X=\{x\in A : \chi(x)=1\}$), so the
    number of indicator functions equals $|\mathcal P(A)| = 2^n$.
\end{solution}

% ======================================================================
\begin{exercise}{2}
    Let $A$ and $B$ be any sets. Express the indicator functions of
    each of the following in terms of the indicator functions
    $\chi_A$ and $\chi_B$.
    \begin{enumerate}[label=(\alph*)]
        \item $\overline{A}$
        \item $A\cap B$
        \item $A\setminus B$
        \item $A\triangle B$
    \end{enumerate}
\end{exercise}
\begin{solution}
    \textbf{(a)}
    \[
        \chi_{\overline A}(x) = 1-\chi_A(x).
    \]

    \textbf{(b)} $x\in A\cap B \iff x\in A \wedge x\in B$, so
    \[
        \chi_{A\cap B} = \chi_A\cdot\chi_B.
    \]

    \textbf{(c)} $x\in A\setminus B \iff x\in A \wedge x\notin B$, so
    \[
        \chi_{A\setminus B} = \chi_A\cdot\chi_{\overline B}
        = \chi_A\cdot(1-\chi_B) = \chi_A - \chi_A\chi_B.
    \]

    \textbf{(d)} Since $A\triangle B = (A\setminus B)\cup(B\setminus
    A)$ is a disjoint union,
    \[
        \chi_{A\triangle B} = \chi_{A\setminus B}+\chi_{B\setminus A}
        = (\chi_A-\chi_A\chi_B)+(\chi_B-\chi_A\chi_B)
        = \chi_A+\chi_B-2\chi_A\chi_B.
    \]
\end{solution}

\begin{exercise}{3}
    If two functions $f:A\to\R$ and $g:A\to\R$ have the same domain
    $A$ and give real numbers as their values, then their sum
    $f+g:A\to\R$ is defined for all $x\in A$ by
    \[
        (f+g)(x) = f(x)+g(x).
    \]
    What is the sum of all the indicator functions of all the
    subsets of a finite set?
\end{exercise}
\begin{proof}
    Let $A$ be a finite set with $|A|=n$. Since function addition is
    defined pointwise, for any $x_i\in A$,
    \[
        \left(\sum_{X\subseteq A}\chi_X\right)(x_i)
        = \sum_{X\subseteq A}\chi_X(x_i).
    \]

    Since $\chi_X(x_i)=1$ if $x_i\in X$ and $0$ otherwise, this sum
    counts the number of subsets $X\subseteq A$ containing $x_i$.

    \textbf{Method 1 (direct count).} A subset containing $x_i$ is
    formed by fixing $x_i$ to be included, and choosing freely,
    independently, whether each of the remaining $n-1$ elements is
    included. This gives $2^{n-1}$ such subsets.

    \textbf{Method 2 (summing over subset sizes).} For each size
    $k=1,\dots,n$, the number of size-$k$ subsets containing $x_i$ is
    obtained by fixing $x_i$ as one member and choosing the
    remaining $k-1$ members from the other $n-1$ elements:
    $\binom{n-1}{k-1}$. Summing over all sizes,
    \[
        \sum_{X\subseteq A}\chi_X(x_i)
        = \sum_{k=1}^{n}\binom{n-1}{k-1}
        = \sum_{j=0}^{n-1}\binom{n-1}{j} = 2^{n-1},
    \]
    using the binomial identity $\sum_{j=0}^{m}\binom{m}{j}=2^m$
    with $m=n-1$.

    Both methods agree: $\left(\sum_{X\subseteq A}\chi_X\right)(x_i)
    = 2^{n-1}$ for every $x_i\in A$, independent of which $x_i$ is
    chosen. So $\sum_{X\subseteq A}\chi_X$ is the constant function
    on $A$ with value $2^{n-1}$.
\end{proof}

% ======================================================================
\begin{exercise}{5}
    Functions can be viewed as sets of ordered pairs, so we can
    combine functions using set operations. The result will be a
    binary relation but, depending on the operation, it may or may
    not be another function.

    Suppose $f: A\to B$ and $g: C\to D$ are functions. Which of the
    following is always a function?
    \[
        f\cap g;\quad f\cup g;\quad f\setminus g;\quad g\setminus f;
        \quad f\triangle g;\quad f\times g.
    \]
    For each that is always a function, give its definition in the
    usual form, including showing how the rule depends on $f$ and
    $g$. For each that is not necessarily a function, explain why
    this is the case (e.g., with the help of examples for $f$ and
    $g$ under which the operation does not give a function).
\end{exercise}
\begin{proof}
    Recall that $R\subseteq A\times B$ is a function if and only if
    $R$ satisfies both:
    \begin{itemize}
        \item \textbf{Existence}: $\{x\in A : (x,y)\in R \text{ for
            some } y\in B\} = A$;
        \item \textbf{Uniqueness (well-definedness)}: for every
            $x\in A$, the set $R(x)=\{y\in B : xRy\}$ contains
            exactly one element.
    \end{itemize}

    \medskip
    \textbf{(i) $f\cap g$.}
    \[
        (f\cap g)(x) = \{f(x)\}\cap\{g(x)\}.
    \]
    When $f(x)=g(x)$: $(f\cap g)(x) = \{f(x)\}$.
    When $f(x)\ne g(x)$: $(f\cap g)(x) = \emptyset$ --- does not
    exist. So $f\cap g$ is \emph{not necessarily} a function; it can
    fail \textbf{existence}.

    \medskip
    \textbf{(ii) $f\cup g$.} \label{5ii}
    \[
        (f\cup g)(x) = \{f(x)\}\cup\{g(x)\}.
    \]
    When $f(x)=g(x)$: $|(f\cup g)(x)|=1$.
    When $f(x)\ne g(x)$: $|(f\cup g)(x)|=2$. So $f\cup g$ is
    \emph{not necessarily} a function; it can fail
    \textbf{uniqueness}.

    \medskip
    \textbf{(iii) $f\setminus g$.}
    \[
        (f\setminus g)(x) = \{f(x)\}\setminus\{g(x)\}.
    \]
    When $f(x)=g(x)$: $|f\setminus g|=0$ --- does not exist.
    When $f(x)\ne g(x)$: $|f\setminus g|=1$. So $f\setminus g$ is
    \emph{not necessarily} a function; it can fail
    \textbf{existence}.

    \medskip
    \textbf{(iv) $g\setminus f$.}
    By the same argument as (iii) with the roles of $f,g$
    interchanged, $g\setminus f$ is \emph{not necessarily} a
    function.

    \medskip
    \textbf{(v) $f\triangle g$.}
    \[
        (f\triangle g)(x) = \{(f\setminus g)(x)\}
        \cup \{(g\setminus f)(x)\}.
    \]
    When $f(x)=g(x)$: $|f\triangle g|=0$ --- does not exist.
    When $f(x)\ne g(x)$: $|f\triangle g|=2$. So $f\triangle g$ is
    \emph{not necessarily} a function; depending on $f,g$, it can
    fail either \textbf{existence} or \textbf{uniqueness}.

    \medskip
    \textbf{(vi) $f\times g$.}
    \[
        f\times g = \{((a,b),(c,d)) : (a,b)\in f,\ (c,d)\in g\}.
    \]
    This defines a function
    \[
        f\times g : A\times C \to B\times D, \qquad
        (f\times g)(x,y) = (f(x), g(y)).
    \]
    \textbf{Existence}: for all $(x,y)\in A\times C$, since $f$ and
    $g$ are functions, $f(x)$ and $g(y)$ must exist, so
    $(f(x),g(y))$ exists.
    \textbf{Uniqueness}: since $f(x)$ and $g(y)$ are each unique,
    the pair $(f(x),g(y))$ is unique.
    So $f\times g$ is \emph{always} a function.
\end{proof}

% ======================================================================
\begin{exercise}{6}
    Suppose that $f: A\to B$ and $g: C\to D$ are functions with
    disjoint domains: $A\cap C=\emptyset$. Is the disjoint union
    $f\sqcup g$ always a function? If so, give its definition in
    terms of $f$ and $g$. If not, why not?
\end{exercise}
\begin{solution}
    By \hyperref[5ii]{5(ii)} $f \cup g $ is always a function.
\end{solution}

% ======================================================================
\begin{exercise}{8}
    A bitstring is a string over the binary alphabet $\{0,1\}$; in
    other words, it is a member of $\{0,1\}^*$.
    \begin{enumerate}[label=(\alph*)]
        \item Give a bijection from $\{0,1\}^*$ to $\N$.

        Now let $\{0,1\}^{*+}$ denote the set of all finite nonempty
        tuples (or sequences) of bitstrings. Examples of members of
        $\{0,1\}^{*+}$ include:
        \[
            (011,1,10001),\quad (10,10,10,10),\quad (\varepsilon,
            1011),\quad (1010010001).
        \]
        \item (Challenge) Give a bijection from $\{0,1\}^{*+}$ to
        $\{0,1\}^*$.
    \end{enumerate}
\end{exercise}
\begin{solution}[(a)]

\end{solution}
\begin{solution}[(b)]

\end{solution}

% ======================================================================
\begin{exercise}{9}
    Let $A$ be a set of $n$ elements and let $x\in A$. We say that a
    function $f: A\to A$ \emph{fixes} $x$ if $f(x)=x$. We also say in
    this case that $x$ is a \emph{fixed point} of $f$.

    This says nothing about what $f$ does to other elements of $A$;
    it is possible that some other elements are fixed by $f$ too, or
    that none are. All this definition requires is that $f$ sends
    $x$ to itself.

    Now let $X\subseteq A$. The function $f$ \emph{fixes} $X$ if it
    fixes all elements of $X$. So, for all $x\in X$, we have
    $f(x)=x$. In other words, the restriction of $f$ to $X$ is just
    the identity function on $X$:
    \[
        f|_X = i_X.
    \]
    Members of $A\setminus X$ may or may not be fixed by $f$.

    Let $\mathrm{Fix}(X)$ be the set of bijections on $A$ that fix
    $X$, and let $F$ be the set of all bijections on $A$ that fix at
    least one element of $A$.

    Using Exercise 1.13 (i.e., Exercise 13 in Chapter 1), express
    $|F|$ in terms of the set sizes $|\mathrm{Fix}(X)|$ using all
    $X\subseteq A$.
\end{exercise}
\begin{solution}

\end{solution}

% ======================================================================
\begin{exercise}{10}
    For each $k\in\{1,\dots,n\}$ and any subset $X\subseteq A$ with
    $|X|=k$, write $f_k$ for the number of bijections on $A$ that fix
    $X$.
    \begin{enumerate}[label=(\alph*)]
        \item Why don't we include $X$ in the notation $f_k$, given
            that its definition refers to $X$?
        \item Give an expression for $f_k$.
        \item Express $|F|$ in terms of $f_1,f_2,\dots,f_n$, and then
            use your expression for $f_k$ to give an expression for
            $|F|$.
        \item Hence express the number of fixed-point-free
            bijections on $A$ in terms of $f_1,f_2,\dots,f_n$. A
            function is fixed-point-free if it has no fixed points.
        \item For each of the following sets, what proportion of all
            bijections on the set are fixed-point-free?
            $\{0,1\}$; $\{\spadesuit,\clubsuit,\diamondsuit,
            \heartsuit\}$; the set of ten decimal digits; the
            26-letter English alphabet.
        \item What is the largest set for which you can determine
            this proportion, and what is the value of the proportion
            for that set?
    \end{enumerate}
\end{exercise}
\begin{solution}[(a)]

\end{solution}
\begin{solution}[(b)]

\end{solution}
\begin{solution}[(c)]

\end{solution}
\begin{solution}[(d)]

\end{solution}
\begin{solution}[(e)]

\end{solution}
\begin{solution}[(f)]

\end{solution}

% ======================================================================
\begin{exercise}{11}
    \begin{enumerate}[label=(\alph*)]
        \item Write down all surjections from $\{1,2,3,4\}$ to
            $\{a,b,c\}$. Compare your answer with Exercise 1.20.
        \item Define a surjection from
        \[
            \{\text{surjections from } \{1,2,3,4\} \text{ to }
            \{a,b,c\}\}
        \]
        to
        \[
            \{\text{partitions of } \{1,2,3,4\} \text{ into three
            parts}\}.
        \]
        \item What sizes can the preimages of members of its image
            be?
    \end{enumerate}
\end{exercise}
\begin{solution}

\end{solution}

% ======================================================================
\begin{exercise}{12}
    Suppose $f:A\to B$ and $g:C\to D$ are bijections, $A\cap
    C=\emptyset$ and $B\cap D=\emptyset$. Since $A\cap C=\emptyset$,
    the disjoint union $f\sqcup g$ exists. Show that $f\sqcup g$ is a
    bijection and express the inverse of $f\sqcup g$ in terms of the
    inverses of $f$ and $g$.
\end{exercise}
\begin{solution}

\end{solution}

% ======================================================================
\begin{exercise}{13}
    If the function $f$ has an inverse function, which of its
    iterated compositions $f^{(k)}$ have inverse functions (where
    $k\in\N$)?
\end{exercise}
\begin{solution}

\end{solution}

% ======================================================================
\begin{exercise}{14}
    Determine all functions from $\{1,2,3\}$ to itself whose
    indefinitely iterated composition is a constant function (i.e.,
    $f^{(k)}$ is a constant function if $k$ is large enough).

    For a challenge: investigate when this happens in general. Try
    to characterise those functions $f:\{1,2,\dots,n\}\to
    \{1,2,\dots,n\}$ such that, for large enough $k$, the iterated
    composition $f^{(k)}$ is constant.
\end{exercise}
\begin{solution}

\end{solution}

% ======================================================================
\begin{exercise}{15}
    This exercise is about the Caesar slide cryptosystem, one of the
    oldest and simplest cryptosystems. It is not secure, but its core
    operation is used in many stronger and more complex
    cryptosystems.

    Caesar slide encrypts a message by sliding each letter along the
    alphabet, rightwards, by some fixed number of steps. This fixed
    number is the key, $k$, and the same amount of sliding is done to
    each letter of the message, with wrap-around.

    Formal definition: the message space $M$ is the set of all
    strings of English lower case letters. The key space $K$ is the
    English lower-case alphabet $\{a,\dots,z\}$. The cypher space $C$
    equals $M$. Letter addition $\alpha+\beta$ slides $\alpha$
    rightwards by $\beta$ steps, with wrap-around (and letter
    subtraction slides leftwards).

    The encryption function $e: M\times K\to M$ is defined for
    $m=m_1m_2\cdots m_n\in M$ and $k\in K$ by
    \[
        e(m,k) = c_1c_2\cdots c_n, \qquad c_i = m_i+k.
    \]
    The decryption function $d: M\times K\to M$ is defined for
    $c=c_1c_2\cdots c_n\in M$ and $k\in K$ by
    \[
        d(c,k) = m_1m_2\cdots m_n, \qquad m_i = c_i-k.
    \]

    Referring to the definitions of \emph{embeddable}, \emph{
    equivalent} and \emph{idempotent} given in Assignment 1: prove
    that the Caesar slide cryptosystem is idempotent.
\end{exercise}
\begin{solution}

\end{solution}

% ======================================================================
\begin{exercise}{16}
    If $f$ is an injection and $f^{-1}$ is its inverse relation, what
    can you say about $f\circ f^{-1}$ and $f^{-1}\circ f$? What kinds
    of functions or relations are they, and what are their domains
    and codomains?
\end{exercise}
\begin{solution}

\end{solution}

% ======================================================================
\begin{exercise}{18}
    Recall the Parent relation from \S2.13. Suggest a good name for
    its transitive closure. Then use that relation (i.e., the
    transitive closure) to define formally the relation that links
    two people that are related to each other (however distantly).
\end{exercise}
\begin{solution}

\end{solution}

% ======================================================================
\begin{exercise}{19}
    Rephrase the Collatz Conjecture as a statement about the
    transitive closure of the Collatz function.
\end{exercise}
\begin{solution}

\end{solution}

% ======================================================================
\begin{exercise}{20}
    Below are some binary relations on the set of all Python
    programs, denoted $P,Q$, with relations $\simeq_1,\dots,
    \simeq_{10}$. For each relation, determine whether or not it is
    an equivalence relation, and give reasons.
    \begin{center}
    \begin{tabular}{ll}
        $P\simeq_1 Q$ & $P$ and $Q$ have the same number of
            characters \\
        $P\simeq_2 Q$ & $P$ and $Q$ compute the same function \\
        $P\simeq_3 Q$ & $P$ and $Q$ have the same set of variable
            names \\
        $P\simeq_4 Q$ & $P$ and $Q$ have at least one variable name
            in common \\
        $P\simeq_5 Q$ & $P$ and $Q$ have no variable names in
            common \\
        $P\simeq_6 Q$ & $P$ and $Q$ were written by the same set of
            programmers \\
        $P\simeq_7 Q$ & $P$ and $Q$ have at least one coauthor in
            common \\
        $P\simeq_8 Q$ & $P$ was completed before $Q$ \\
        $P\simeq_9 Q$ & $P$ was completed on the same day as $Q$ \\
        $P\simeq_{10} Q$ & $P$ was completed after $Q$ \\
    \end{tabular}
    \end{center}
\end{exercise}
\begin{solution}

\end{solution}

% ======================================================================
\begin{exercise}{22}
    The anagram relation on the set of all English words is defined
    as follows. Let $x_1x_2\cdots x_m$ and $y_1y_2\cdots y_n$ be two
    English words of lengths $m$ and $n$ respectively. The ordered
    pair $(x_1\cdots x_m, y_1\cdots y_n)$ belongs to the anagram
    relation if and only if there exists a bijection
    $f:\{1,\dots,m\}\to\{1,\dots,n\}$ such that, for all
    $i\in\{1,\dots,m\}$, $x_i=y_{f(i)}$.
    \begin{enumerate}[label=(\alph*)]
        \item Find the largest set of three-letter words you can
            that are all related to each other by the anagram
            relation.
        \item For each $k=1,2,3,4$, find an English word of length
            $k$ that is not related to any other word by anagram.
        \item Prove that anagram is an equivalence relation.
    \end{enumerate}
\end{exercise}
\begin{solution}[(a)]

\end{solution}
\begin{solution}[(b)]

\end{solution}
\begin{solution}[(c)]

\end{solution}

% ======================================================================
\begin{exercise}{24}
    Recall that when you have a table with $n$ columns, its rows
    represent the $n$-tuples in an $n$-ary relation. For $1\le i\le
    n$, let $A_i$ be a set containing all the entries in column $i$
    (and possibly more elements), so that the columns of the table
    define an $n$-ary relation on $A_1\times A_2\times\cdots\times
    A_n$. Assuming the columns of this table are all distinct, how
    many ternary relations can you construct by choosing columns from
    this table of $n$ columns?
\end{exercise}
\begin{solution}

\end{solution}

\end{document}